\documentclass[conference]{IEEEtran}
\IEEEoverridecommandlockouts

\usepackage{cite}
\usepackage{amsmath,amssymb,amsfonts}
\usepackage{algorithmic}
\usepackage{graphicx}
\usepackage{textcomp}
\usepackage{xcolor}
\usepackage{booktabs}
\usepackage{multirow}
\usepackage{url}
\usepackage{microtype}

\def\BibTeX{{\rm B\kern-.05em{\sc i\kern-.025em b}\kern-.08em
    T\kern-.1667em\lower.7ex\hbox{E}\kern-.125emX}}

\begin{document}

\title{InfoShield: From Graph-Based Misinformation Detection\\
to Community-Aware Cascade Suppression}

\author{
  \IEEEauthorblockN{[Student Name 1]\IEEEauthorrefmark{1},
                    [Student Name 2]\IEEEauthorrefmark{1},
                    [Student Name 3]\IEEEauthorrefmark{1}}
  \IEEEauthorblockA{\IEEEauthorrefmark{1}[Department / Institution Placeholder]\\
  [City, Country]\\
  \{[email1], [email2], [email3]\}@[institution.edu]}
}

\IEEEtitleabstractindextext{%
\begin{abstract}
\centering
Misinformation on social media spreads through community-specific cascade
structures that differ enough between platforms to invalidate models trained on
one community and applied to another.
InfoShield is a four-phase pipeline that addresses this by combining a
Bi-directional Graph Convolutional Network (BiGCN) for cascade classification,
a Stochastic Block Model (SBM) for community-specific spreading dynamics, and
a Linear Programming (LP) optimizer that computes per-edge dropout policies to
suppress false-content spread while preserving true-content propagation.
Phase~1 establishes that Twitter-15 and Weibo differ by 71\% in mean cascade
width and 31\% in depth (\(p < 0.001\), Mann-Whitney~U), confirming that
community-specific parameters are statistically necessary.
Within-Weibo, rumour and non-rumour cascades differ by a factor of 4.1$\times$
in depth-to-width ratio---a structural signal exploitable by the GNN.
Phase~2 achieves 0.83 accuracy and 0.76 macro-F1 on partial 5-fold
cross-validation on Twitter-15, improving over the previous version by
\(+0.06\) accuracy and \(+0.08\) F1.
Phases~3 and~4 are implemented but not yet evaluated end-to-end.
\end{abstract}

\begin{IEEEkeywords}
misinformation detection, graph convolutional networks, stochastic block model,
cascade suppression, linear programming, social media
\end{IEEEkeywords}}

\maketitle
\IEEEdisplaynontitleabstractindextext
\IEEEpeerreviewmaketitle

%% ---------------------------------------------------------------
\section{Introduction \& Motivation}

The engineering problem is not simply detecting misinformation---it is that
spreading dynamics differ structurally between communities, so a model
calibrated on one community's cascade priors will be miscalibrated when
deployed on another.
Concretely: Twitter-15 cascades have a mean width of 309.84 and mean depth of
3.96, while Weibo cascades have a mean width of 530.53 and mean depth of 5.19.
That is a 71\% difference in width and a 31\% difference in depth.
A suppression policy derived from Twitter-15's SBM parameters would
underestimate the reach of false content on Weibo by the same margin before
a single node feature is examined.
These differences are not sampling artefacts: Mann-Whitney~U and
Kruskal-Wallis tests both return \(p < 0.001\) for all three primary metrics
(depth, width, depth/width ratio).

InfoShield addresses this through three stacked layers:

\begin{itemize}
  \item \textbf{Detection layer (BiGCN)}: Given a propagation tree rooted at a
    source tweet, classify the cascade as \textit{true}, \textit{false},
    \textit{unverified}, or \textit{non-rumour}.
  \item \textbf{Network model layer (SBM)}: Given GNN-labelled cascades, fit
    community-specific branching matrices \(b^+\) (true content) and \(b^-\)
    (false content) over \(k\) polarization classes.
  \item \textbf{Intervention layer (LP)}: At each SIR time step, solve an LP
    that selects a per-edge dropout matrix \(d^*\) to minimize false-content
    spread subject to a minimum branching ratio for true content.
\end{itemize}

Each layer must be community-specific because the SBM transition rates and
the GNN node-feature distributions are platform-dependent.
A single global matrix would average out the structural differences that
Phase~1 quantifies, making the LP optimizer solve the wrong problem.

The target evaluation community is WICO---COVID-19 conspiracy tweets on
English-language Twitter.
Twitter-15 (English general rumours) serves as the training domain for the
BiGCN.
Weibo (Chinese Sina Weibo) provides a structurally distinct third community
for cross-community comparison.

%% ---------------------------------------------------------------
\section{Related Work}

\textbf{BiGCN (Bian et al.~\cite{bian2020rumor})}.
The original BiGCN propagates representations both top-down and bottom-up
through a cascade tree, then concatenates the pooled embeddings for
classification.
We adapt this architecture by (a)~appending four structural node features to
the 768-dim RoBERTa CLS embedding and (b)~injecting the root representation at
Layer~2 of each branch, extending effective context depth beyond the 2-hop GCN
limit.
The classification head and training protocol are otherwise consistent with
the original.

\textbf{SBM cascade models}.
Stochastic block models of information diffusion partition users into
structurally homogeneous blocks and estimate block-to-block transfer rates.
This gives a compact, interpretable representation of spreading dynamics that
separates by content label---the key property we exploit when fitting separate
\(b^+\) and \(b^-\) matrices.

\textbf{LP-based intervention}.
Optimal edge-dropout under spreading-rate constraints is a linear program when
the SIR transition rates are fixed.
We follow the formulation of Algorithm~2 in~\cite{infoshield_algo}, interleaving
LP solves with the SIR simulation so the dropout policy adapts to the current
infected/susceptible population at each time step.

%% ---------------------------------------------------------------
\section{System Architecture}

The pipeline is a directed acyclic computation graph:

\begin{enumerate}
  \item Raw propagation trees (one file per cascade) are parsed into
    PyTorch-Geometric \texttt{Data} objects with 772-dim node features.
  \item The BiGCN classifier assigns a 4-class label and confidence score to
    each cascade.
  \item Cascades with confidence \(\geq 0.65\) feed the SBM fitter, which
    builds the union WICO graph, runs Louvain clustering (\texttt{resolution=2.0},
    \(k=13\)), and estimates \(b^+\) and \(b^-\) via frequentist MLE.
  \item At each SIR time step, the LP optimizer receives \(b^+\), \(b^-\), and
    the current compartment counts \(|S^v_t|\), \(|I^u_t|\), and solves for
    \(d^*\).
  \item \(d^*\) is applied as an edge mask to the SIR update equation.
\end{enumerate}

The bidirectional dependency is explicit: the LP needs \(b^+\) and \(b^-\)
from Phase~3, which requires Phase~2 predictions to split cascade edges by
label before fitting.
No phase can be short-circuited.

%% ---------------------------------------------------------------
\section{Phase 1: Cross-Community Cascade Analysis}

\subsection{Twitter-15 vs.\ Weibo}

Table~\ref{tab:cross_community} reports the four structural metrics for
Twitter-15 and Weibo.

\begin{table}[h]
\caption{Cross-community cascade structure comparison.}
\label{tab:cross_community}
\centering
\small
\begin{tabular}{lrr}
\toprule
\textbf{Metric} & \textbf{Twitter-15} & \textbf{Weibo} \\
\midrule
Mean depth            & 3.96  & 5.19  \\
Mean width            & 309.84 & 530.53 \\
Depth/width ratio     & 0.023 & 0.054 \\
Gini (depth)          & 0.21  & 0.35  \\
Gini (width)          & 0.454 & 0.715 \\
Mean cascade size     & 402.15 & 790.30 \\
Dataset size          & 1,490 & 4,659 \\
\bottomrule
\end{tabular}
\end{table}

All three primary metrics differ with \(p < 0.001\) under both Mann-Whitney~U
and Kruskal-Wallis tests.
The Gini coefficients compound this: Weibo's width distribution (\(G=0.715\))
is far more concentrated than Twitter-15's (\(G=0.454\)), meaning a small
fraction of Weibo cascades account for most of the total spread.
A suppression policy that treats all cascades equally would misallocate
intervention budget on Weibo.

\subsection{Within-Weibo Label Analysis}

Table~\ref{tab:weibo_labels} compares rumour and non-rumour cascades within
Weibo.

\begin{table}[h]
\caption{Within-Weibo structural comparison by label.}
\label{tab:weibo_labels}
\centering
\small
\begin{tabular}{lrr}
\toprule
\textbf{Metric} & \textbf{Rumour} & \textbf{Non-Rumour} \\
\midrule
Mean depth         & 3.91  & 6.72  \\
Mean width         & 566.6 & 494.7 \\
Depth/width ratio  & 0.022 & 0.090 \\
Mean size          & 721.2 & 876.2 \\
\bottomrule
\end{tabular}
\end{table}

The depth/width ratio separates the two classes by 4.1$\times$ (0.090 vs.\
0.022).
Non-rumour cascades are deeper relative to their width---they propagate through
longer chains rather than broad fan-out.
Rumour cascades spread wide quickly, then stall.
This structural signal is available to the GNN via the four per-node structural
features, independent of text content.

Weibo label balance: 2,278 rumour (48.9\%), 2,249 non-rumour (48.3\%),
132 unlabelled (2.8\%).
The dataset is near-balanced, so accuracy and macro-F1 are comparable metrics.

\subsection{Limitations of Phase 1}

Temporal dynamics (spread rate over time) are not yet analysed.
WICO structural metrics are not yet computed.
Missing-label bias in the 2.8\% unlabelled Weibo fraction is not characterised.

%% ---------------------------------------------------------------
\section{Phase 2: BiGCN Classification}

\subsection{Architecture}

Each cascade is represented as a directed tree graph.
Node features are 772-dimensional: 768 dimensions from a frozen RoBERTa-base
CLS token (broadcast from root to all nodes), concatenated with four structural
features per node---\texttt{is\_root}, \texttt{norm\_depth},
\texttt{norm\_breadth}, and \texttt{norm\_descendants}.

Two parallel \texttt{GCNBranch} modules process top-down and bottom-up
edge orientations separately:

\begin{enumerate}
  \item \textbf{Layer 1}: \texttt{GCNConv}(772 $\to$ 256) + ReLU + Dropout(0.5)
  \item \textbf{Root augmentation}: root features (772-dim) concatenated to
    every node $\Rightarrow$ 1028-dim per node
  \item \textbf{Layer 2}: \texttt{GCNConv}(1028 $\to$ 128) + ReLU + Dropout(0.5)
  \item \textbf{Pooling}: global mean pool $\Rightarrow$ (batch, 128)
\end{enumerate}

The two pooled representations are concatenated to (batch, 256), then passed
through \texttt{Linear}(256$\to$128) + ReLU + Dropout(0.5) and
\texttt{Linear}(128$\to$4) for logit output.

\textbf{Root enhancement rationale.}
A standard 2-layer GCN propagates information at most 2 hops from any node.
With mean cascade depths of 3.96--5.19, nodes beyond hop 2 from the source
lose source-tweet context after Layer~1.
Concatenating the 772-dim root embedding to every node's intermediate
representation at Layer~2 re-injects source context globally, so the
bidirectional branches can modulate deep subtrees with the original post's
semantics.
Without this augmentation, the classifier relies on propagated structural
signals alone for deep nodes.

\textbf{Frozen RoBERTa.}
We froze RoBERTa-base (\(\approx\)125M parameters) rather than fine-tuning it.
Fine-tuning on 1,490 cascades would overfit: the ratio of trainable parameters
to training examples would be roughly 84,000:1, which no regularisation scheme
corrects adequately.
Frozen CLS embeddings provide consistent, pre-trained semantic representations.
The GNN layers learn to discriminate on top of them.

\subsection{Training Setup}

Table~\ref{tab:hyperparams} lists the hyperparameters and their explicit
justifications.

\begin{table}[h]
\caption{Hyperparameter settings and justifications.}
\label{tab:hyperparams}
\centering
\small
\begin{tabular}{lll}
\toprule
\textbf{Param} & \textbf{Value} & \textbf{Reason} \\
\midrule
lr              & 5e-4 & Standard Adam, frozen encoder \\
weight\_decay   & 5e-4 & L2 reg., small dataset \\
batch\_size     & 32   & CPU memory; acceptable noise \\
dropout         & 0.5  & Aggressive; standard for small GNNs \\
patience        & 20   & Prevent overfit, small data \\
label\_smoothing& 0.05 & Noisy social media labels \\
hidden\_dim     & 256  & Capacity vs.\ overfit balance \\
out\_dim        & 128  & Standard for cascade classification \\
\bottomrule
\end{tabular}
\end{table}

Training uses 5-fold CV on Twitter-15.
Each fold trains with Adam, cross-entropy loss with label smoothing 0.05, and
early stopping at patience 20.
Per-fold checkpoints are saved at the best validation loss.

\subsection{Results}

Table~\ref{tab:bigcn_results} reports Phase~2 results.

\begin{table}[h]
\caption{BiGCN classification results on Twitter-15 (partial 5-fold CV).}
\label{tab:bigcn_results}
\centering
\small
\begin{tabular}{lr}
\toprule
\textbf{Metric} & \textbf{Value} \\
\midrule
4-class accuracy    & 0.83 \\
4-class macro F1    & 0.76 \\
Folds completed     & 2 of 5 (fold 0, fold 1) \\
Previous accuracy   & 0.77 \\
Previous macro F1   & 0.68 \\
Improvement         & $+0.06$ acc, $+0.08$ F1 \\
Training device     & CPU \\
\bottomrule
\end{tabular}
\end{table}

Folds~0 and~1 completed; folds~2--4 were not run due to wall-clock budget on
CPU.
The 0.83/0.76 figures are therefore not final cross-validated estimates.
The published BiGCN baseline (Bian et al.) exceeds these numbers; the gap is
attributed to CPU-only training limiting the epoch budget before early stopping,
not to a fundamental architectural deficiency.

\subsection{Limitations}

Full 5-fold CV is not complete.
No hyperparameter search beyond the defaults in Table~\ref{tab:hyperparams}
has been performed.
Weibo has not been integrated into BiGCN training; it is used in Phase~1
analysis only.
GPU access would resolve the wall-clock bottleneck.

%% ---------------------------------------------------------------
\section{Phase 3: Stochastic Block Model}

\subsection{Formulation}

The SBM partitions WICO users into \(k=13\) polarization classes via Louvain
modularity clustering on the union graph of all WICO cascade edges
(\texttt{resolution=2.0}).
Partitions containing fewer than 1\% of total users are merged into the nearest
class to prevent near-empty blocks from producing unstable rate estimates.

Two \(k \times k\) transition matrices are fitted:
\[
  b^+_{uv} = \frac{\text{transfers from class-}u\text{ to class-}v
                   \text{ in true cascades}}{|C_u|}
\]
\[
  b^-_{uv} = \frac{\text{transfers from class-}u\text{ to class-}v
                   \text{ in false cascades}}{|C_u|}
\]
This is frequentist MLE: the denominator \(|C_u|\) normalises by the total
number of users in class \(u\), giving the empirical conditional transfer rate.

Only cascades with GNN confidence \(\geq 0.65\) feed the fitter.
Below this threshold, the label is too uncertain to cleanly separate
\(b^+\) and \(b^-\) edges.
Setting the threshold at 0.65 rather than 0.5 accepts a recall penalty on
ambiguous cascades in exchange for cleaner matrix estimates.

\subsection{Status}

The SBM fitter is implemented and unit-tested on synthetic data.
It has not yet been run on the full WICO graph.
No numerical results are available for \(b^+\) or \(b^-\).

%% ---------------------------------------------------------------
\section{Phase 4: LP-Based Dropout Optimization}

\subsection{Algorithm 2}

Algorithm~2 interleaves LP solves with the discrete SIR simulation.
At time step \(t\), given current susceptible counts \(|S^v_t|\) and infected
counts \(|I^u_t|\) for each class pair \((u,v)\), two LPs are available.

\textbf{Primary LP} (minimize false spread, preserve true spread):
\begin{align}
  \min_{d} \quad & \sum_{u,v} |S^v_t|\,|I^u_t|\,d_{uv}\,b^-_{uv} \label{eq:primary_obj} \\
  \text{s.t.} \quad
    & \sum_{u,v} |S^v_t|\,|I^u_t|\,d_{uv}\,b^+_{uv} \geq \alpha|I_t| \label{eq:primary_con} \\
    & d_{uv} \in [0,1]^{k \times k} \notag
\end{align}

\textbf{Softened LP} (fallback when primary is infeasible):
\begin{align}
  \min_{d} \quad & \sum_{u,v} |S^v_t|\,|I^u_t|\bigl(b^-_{uv} + \lambda b^+_{uv}\bigr)d_{uv}
  \label{eq:soft_obj}
\end{align}
with the same box constraints on \(d\).

The primary LP becomes infeasible when true content is already spreading slowly
enough that no dropout policy can satisfy the \(\alpha\)-constraint---i.e.,
the system cannot simultaneously suppress false content and guarantee the
required true-content branching ratio.
In that regime, the softened LP minimizes a weighted sum of false suppression
and true preservation with equal weight \(\lambda=1.0\).

\textbf{Parameter choices.}
\(\alpha=1.5\) requires the selected dropout policy to sustain a branching ratio
of at least 1.5 for true content, preventing the optimizer from achieving false
suppression by simply muting all edges.
\(\lambda=1.0\) assigns equal penalty per unit of true content reduced and per
unit of false content passed---a neutral prior that can be tuned once
suppression data is available.
Both are configurable.

Both LPs are solved with \texttt{scipy.optimize.linprog} using the
interior-point method.

\subsection{Status}

The LP optimizer is implemented.
End-to-end Algorithm~2 evaluation on WICO has not been run.
No suppression effectiveness measurements are available.

%% ---------------------------------------------------------------
\section{Engineering Decisions}

\begin{itemize}

  \item \textbf{Frozen vs.\ fine-tuned RoBERTa.}
    Decision: frozen.
    Alternative: end-to-end fine-tuning.
    Reason: 1,490 training cascades cannot constrain a 125M-parameter encoder.
    Fine-tuning would memorise training cascades, not generalise.

  \item \textbf{Root enhancement at Layer 2 vs.\ Layer 1 or omitted.}
    Decision: augment at Layer 2 (after the first GCN hop).
    Alternative: augment at Layer 1 input (no additional hop) or not at all.
    Reason: augmenting at Layer 1 would give every node a direct copy of root
    features, defeating the purpose of learning structural propagation.
    Augmenting at Layer 2 allows Layer 1 to build local structural
    representations, then conditions Layer~2 on global source context.

  \item \textbf{Two-phase LP (primary + softened) vs.\ single LP with slack.}
    Decision: two-phase.
    Alternative: add a slack variable to constraint~\eqref{eq:primary_con} with
    a large penalty.
    Reason: a slack-variable formulation mixes the primary and fallback
    objectives into a single weighted sum, making the tradeoff implicit and
    hard to inspect.
    Two separate LPs make each regime's objective function explicit and allow
    independent parameter tuning.

  \item \textbf{Frequentist MLE for SBM vs.\ Bayesian estimation.}
    Decision: frequentist MLE.
    Alternative: Bayesian with a Dirichlet prior over transfer rates.
    Reason: the WICO graph is large enough that maximum likelihood is not
    variance-dominated.
    Bayesian estimation would add implementation complexity and require
    prior elicitation without a clear empirical benefit at this dataset scale.

  \item \textbf{Louvain resolution 2.0 (\(k=13\)) vs.\ lower resolution.}
    Decision: resolution 2.0, yielding \(k=13\).
    Alternative: resolution 1.0 (default, \(k \approx\) 5--8 on this graph).
    Reason: coarser partition loses within-class heterogeneity in transfer
    rates.
    Resolution 2.0 produces finer-grained classes that better separate high-
    and low-polarization user groups, giving the LP more structured information
    to act on.

  \item \textbf{Community-specific \(b^+/b^-\) vs.\ global matrices.}
    Decision: per-community.
    Alternative: fit a single pair of matrices across all platforms.
    Reason: Phase~1 establishes that Twitter-15 and Weibo differ by 71\% in
    width and 31\% in depth (\(p < 0.001\)).
    A global matrix pools structurally incompatible communities, producing
    estimates that are wrong for each of them.

\end{itemize}

%% ---------------------------------------------------------------
\section{Current Status \& Roadmap}

Table~\ref{tab:status} summarises the completion state of each phase.

\begin{table}[h]
\caption{Phase completion status.}
\label{tab:status}
\centering
\small
\begin{tabular}{p{3.2cm}p{1.5cm}p{2.0cm}}
\toprule
\textbf{Phase} & \textbf{Status} & \textbf{Blocker} \\
\midrule
Ph.1: Cascade analysis (T15 vs.\ Weibo) & Complete & --- \\
Ph.1: WICO structural analysis          & Not done & Deprioritised \\
Ph.2: BiGCN training (2/5 folds)        & Partial  & CPU-only speed \\
Ph.2: Full 5-fold CV                    & Blocked  & Same \\
Ph.2: Weibo BiGCN training              & Not started & --- \\
Ph.3: SBM fitting on WICO               & Not run  & Needs Ph.2 output \\
Ph.4: LP end-to-end evaluation          & Not run  & Needs Ph.3 matrices \\
\bottomrule
\end{tabular}
\end{table}

The critical path is: (a) complete BiGCN folds 2--4 on GPU,
(b) run the SBM fitter on WICO with Phase~2 predictions,
(c) run Algorithm~2 and measure suppression effectiveness---reduction in false
cascade final size---against a no-dropout baseline.
All three steps are blocked only by compute access and sequential data
dependencies, not by unresolved design decisions.

%% ---------------------------------------------------------------
\section{Conclusion}

\textbf{Empirically established:}
\begin{itemize}
  \item Cross-community structural differences are large (71\% width, 31\%
    depth, \(p < 0.001\)) and statistically justify community-specific SBM
    parameters.
  \item Within-Weibo, rumour and non-rumour cascades differ 4.1$\times$ in
    depth/width ratio, confirming that cascade shape is a discriminating feature
    independent of text content.
  \item BiGCN achieves 0.83 accuracy and 0.76 macro-F1 on partial 5-fold CV on
    Twitter-15, with improvements of \(+0.06\) accuracy and \(+0.08\) F1 over
    the previous version.
\end{itemize}

\textbf{Not yet measured:}
\begin{itemize}
  \item SBM fitting quality on WICO (no \(b^+/b^-\) estimates yet).
  \item LP suppression effectiveness: reduction in false-content cascade size.
  \item Tradeoff curve between false suppression and true-content collateral
    damage as \(\alpha\) and \(\lambda\) vary.
  \item System performance on Weibo (Chinese text, no RoBERTa fine-tuning on
    Chinese).
\end{itemize}

The architecture is complete.
The remaining work is execution: finishing cross-validation, running the SBM
fitter, and running Algorithm~2.
None of the outstanding items require new design decisions.

%% ---------------------------------------------------------------
\bibliographystyle{IEEEtran}
\begin{thebibliography}{9}

\bibitem{bian2020rumor}
T.\ Bian, X.\ Xiao, T.\ Xu, P.\ Zhao, W.\ Huang, Y.\ Rong, and J.\ Huang,
``Rumor detection on social media with bi-directional graph convolutional
networks,''
\textit{Proceedings of the AAAI Conference on Artificial Intelligence},
vol.~34, no.~01, pp.~549--556, 2020.

\bibitem{infoshield_algo}
[Algorithm~2 reference -- insert citation for the LP-based SIR suppression
formulation used in this work.]

\end{thebibliography}

\end{document}
